% GENERATED FILE. Edit canonical lessons, then run npm run generate:books.
% canonical-source-hash: fnv1a64:0f5829466d78b860
% canonical-lessons: TA-C425-kalankarai-vilakkam, TA-C425-kolarankam, TA-C425-talam, TA-C425-katci, TA-C425-cirrunticcalai

\chapter{Sights Worth Seeing}
\label{ch:ta-sights-worth-seeing}
\hlchaptermodality{425}

\begin{chapteropening}
\textbf{By the end of this chapter:} \emph{I can talk about a lighthouse, a planetarium, sites, views and a canteen.}

Complete the last lesson of chapter 425: I can talk about a lighthouse, a planetarium, sites, views and a canteen.
\end{chapteropening}

\section[\texorpdfstring{\ta{கலங்கரை} \ta{விளக்கம்} (kalaṅkarai viḷakkam)}{கலங்கரை விளக்கம் (kalaṅkarai viḷakkam)}]{\ta{கலங்கரை} \ta{விளக்கம்} (kalaṅkarai viḷakkam) — a lighthouse}

\label{lesson:TA-C425-kalankarai-vilakkam}



\begin{quote}
Before the new one: say the Tamil for a mansion, then the Tamil for an embassy.
\end{quote}

\subsection*{You'll want to know: \ta{கலங்கரை} \ta{விளக்கம்}}
\textbf{\ta{கலங்கரை} \ta{விளக்கம்}} — \emph{kalaṅkarai viḷakkam} — \textquotedblleft{}a lighthouse\textquotedblright{}.

\begin{quote}
\ta{கலங்கரை} \ta{விளக்கம்} \ta{இரவில்} \ta{கப்பல்களுக்கு} \ta{வழி} \ta{காட்டுகிறது}.
\end{quote}

\begin{quote}
\emph{kalaṅkarai viḷakkam iravil kappalkaḷukku vaḻi kāṭṭukiṟatu.}
\end{quote}

\begin{quote}
The lighthouse shows ships the way at night.
\end{quote}

Two words: \ta{கலங்கரை}, calling ships, + \ta{விளக்கம்}, a light.

\begin{cousinweb}
\textbf{\ta{கலங்கரை} \ta{விளக்கம்}} (\emph{kalaṅkarai viḷakkam}) is usually traced to \ta{கலம்}, \textquotedblleft{}a vessel, a ship\textquotedblright{} (DEDR 1305), + \ta{கரை}, \textquotedblleft{}to call\textquotedblright{} (DEDR 1291), + \ta{விளக்கம்}, \textquotedblleft{}a light\textquotedblright{}: the light that calls ships in.
\end{cousinweb}

\begin{grammarlens}[title={What you've built}]
One more word for getting around.
\end{grammarlens}

\subsection*{Your turn}
\begin{itemize}
\raggedright
  \item \emph{Say it:} \emph{kalaṅkarai viḷakkam}
  \item \emph{Say it:} \emph{kalaṅkarai viḷakkam}, once more
  \item \emph{Say it:} \emph{kalaṅkarai viḷakkam iravil kappalkaḷukku vaḻi kāṭṭukiṟatu.} — The lighthouse shows ships the way at night.
  \item \emph{Recall:} say the Tamil for a mansion, then the Tamil for an embassy, then say \emph{kalaṅkarai viḷakkam} again
\end{itemize}

\begin{tcolorbox}[breakable,colback=teal!4,colframe=teal!35!black,title={Before you move on}]
What does \ta{கலங்கரை} \ta{விளக்கம்} mean? (\textquotedblleft{}A lighthouse\textquotedblright{}.) Say it once more.
\end{tcolorbox}
\section[\texorpdfstring{\ta{கோளரங்கம்} (kōḷaraṅkam)}{கோளரங்கம் (kōḷaraṅkam)}]{\ta{கோளரங்கம்} (kōḷaraṅkam) — a planetarium}

\label{lesson:TA-C425-kolarankam}



\begin{quote}
Before the new one: say the Tamil for an embassy, then the Tamil for a lighthouse.
\end{quote}

\subsection*{You'll want to know: \ta{கோளரங்கம்}}
\textbf{\ta{கோளரங்கம்}} — \emph{kōḷaraṅkam} — \textquotedblleft{}a planetarium\textquotedblright{}.

\begin{quote}
\ta{கோளரங்கத்தில்} \ta{நட்சத்திரங்களைப்} \ta{பற்றிக்} \ta{கற்றுக்கொண்டோம்}.
\end{quote}

\begin{quote}
\emph{kōḷaraṅkattil naṭcattiraṅkaḷaip paṟṟik kaṟṟukkoṇṭōm.}
\end{quote}

\begin{quote}
At the planetarium we learnt about the stars.
\end{quote}

\ta{கோள்}, a planet, + a word for a hall or a stage.

\begin{cousinweb}
\textbf{\ta{கோளரங்கம்}} (\emph{kōḷaraṅkam}) is a modern coinage built from \ta{கோள்}, \textquotedblleft{}a planet\textquotedblright{}, and \ta{அரங்கம்}, \textquotedblleft{}a hall, a stage\textquotedblright{}, borrowed from Sanskrit \emph{raṅga}, \textquotedblleft{}an arena, a stage\textquotedblright{} (CDIAL 10561).
\end{cousinweb}

\begin{grammarlens}[title={What you've built}]
One more word for getting around.
\end{grammarlens}

\subsection*{Your turn}
\begin{itemize}
\raggedright
  \item \emph{Say it:} \emph{kōḷaraṅkam}
  \item \emph{Say it:} \emph{kōḷaraṅkam}, once more
  \item \emph{Say it:} \emph{kōḷaraṅkattil naṭcattiraṅkaḷaip paṟṟik kaṟṟukkoṇṭōm.} — At the planetarium we learnt about the stars.
  \item \emph{Recall:} say the Tamil for an embassy, then the Tamil for a lighthouse, then say \emph{kōḷaraṅkam} again
\end{itemize}

\begin{tcolorbox}[breakable,colback=teal!4,colframe=teal!35!black,title={Before you move on}]
What does \ta{கோளரங்கம்} mean? (\textquotedblleft{}A planetarium\textquotedblright{}.) Say it once more.
\end{tcolorbox}
\section[\texorpdfstring{\ta{தலம்} (talam)}{தலம் (talam)}]{\ta{தலம்} (talam) — a site; a place}

\label{lesson:TA-C425-talam}



\begin{quote}
Before the new one: say the Tamil for a lighthouse, then the Tamil for a planetarium.
\end{quote}

\subsection*{You'll want to know: \ta{தலம்}}
\textbf{\ta{தலம்}} — \emph{talam} — \textquotedblleft{}a site; a place\textquotedblright{}.

\begin{quote}
\ta{இது} \ta{ஒரு} \ta{முக்கியமான} \ta{சுற்றுலாத்} \ta{தலம்}.
\end{quote}

\begin{quote}
\emph{itu oru mukkiyamāṉa cuṟṟulāt talam.}
\end{quote}

\begin{quote}
This is an important tourist site.
\end{quote}

\ta{சுற்றுலாத்} \ta{தலம்} is a tourist spot.

\begin{cousinweb}
\textbf{\ta{தலம்}} (\emph{talam}) is borrowed from Sanskrit \emph{sthala}, \textquotedblleft{}dry land, ground, a place\textquotedblright{} (CDIAL 13744).
\end{cousinweb}

\begin{grammarlens}[title={What you've built}]
One more word for getting around.
\end{grammarlens}

\subsection*{Your turn}
\begin{itemize}
\raggedright
  \item \emph{Say it:} \emph{talam}
  \item \emph{Say it:} \emph{talam}, once more
  \item \emph{Say it:} \emph{itu oru mukkiyamāṉa cuṟṟulāt talam.} — This is an important tourist site.
  \item \emph{Recall:} say the Tamil for a lighthouse, then the Tamil for a planetarium, then say \emph{talam} again
\end{itemize}

\begin{tcolorbox}[breakable,colback=teal!4,colframe=teal!35!black,title={Before you move on}]
What does \ta{தலம்} mean? (\textquotedblleft{}A site; a place\textquotedblright{}.) Say it once more.
\end{tcolorbox}
\section[\texorpdfstring{\ta{காட்சி} (kāṭci)}{காட்சி (kāṭci)}]{\ta{காட்சி} (kāṭci) — a view; a scene; a sight}

\label{lesson:TA-C425-katci}



\begin{quote}
Before the new one: say the Tamil for a planetarium, then the Tamil for a site.
\end{quote}

\subsection*{You'll want to know: \ta{காட்சி}}
\textbf{\ta{காட்சி}} — \emph{kāṭci} — \textquotedblleft{}a view; a scene; a sight\textquotedblright{}.

\begin{quote}
\ta{மலையிலிருந்து} \ta{பார்த்த} \ta{காட்சி} \ta{மிக} \ta{அழகாக} \ta{இருந்தது}.
\end{quote}

\begin{quote}
\emph{malaiyiliruntu pārtta kāṭci mika aḻakāka iruntatu.}
\end{quote}

\begin{quote}
The view from the hill was very beautiful.
\end{quote}

Built on \ta{காண்}, to see; a scene in a film is also a \ta{காட்சி}.

\begin{cousinweb}
\textbf{\ta{காட்சி}} (\emph{kāṭci}) is native Dravidian (DEDR 1443), the noun of \ta{காண்}, \textquotedblleft{}to see\textquotedblright{}.
\end{cousinweb}

\begin{grammarlens}[title={What you've built}]
One more word for getting around.
\end{grammarlens}

\subsection*{Your turn}
\begin{itemize}
\raggedright
  \item \emph{Say it:} \emph{kāṭci}
  \item \emph{Say it:} \emph{kāṭci}, once more
  \item \emph{Say it:} \emph{malaiyiliruntu pārtta kāṭci mika aḻakāka iruntatu.} — The view from the hill was very beautiful.
  \item \emph{Recall:} say the Tamil for a planetarium, then the Tamil for a site, then say \emph{kāṭci} again
\end{itemize}

\begin{tcolorbox}[breakable,colback=teal!4,colframe=teal!35!black,title={Before you move on}]
What does \ta{காட்சி} mean? (\textquotedblleft{}A view; a scene; a sight\textquotedblright{}.) Say it once more.
\end{tcolorbox}
\section[\texorpdfstring{\ta{சிற்றுண்டிச்சாலை} (ciṟṟuṇṭiccālai)}{சிற்றுண்டிச்சாலை (ciṟṟuṇṭiccālai)}]{\ta{சிற்றுண்டிச்சாலை} (ciṟṟuṇṭiccālai) — a canteen; a snack bar}

\label{lesson:TA-C425-cirrunticcalai}



\begin{quote}
Before the new one: say the Tamil for a site, then the Tamil for a view.
\end{quote}

\subsection*{You'll want to know: \ta{சிற்றுண்டிச்சாலை}}
\textbf{\ta{சிற்றுண்டிச்சாலை}} — \emph{ciṟṟuṇṭiccālai} — \textquotedblleft{}a canteen; a snack bar\textquotedblright{}.

\begin{quote}
\ta{நிலையத்தின்} \ta{சிற்றுண்டிச்சாலையில்} \ta{வடை} \ta{வாங்கினோம்}.
\end{quote}

\begin{quote}
\emph{nilaiyattiṉ ciṟṟuṇṭiccālaiyil vaṭai vāṅkiṉōm.}
\end{quote}

\begin{quote}
We bought vadai at the station canteen.
\end{quote}

\ta{சிறு}, small, + \ta{உண்டி}, food, + \ta{சாலை}, a hall.

\begin{cousinweb}
\textbf{\ta{சிற்றுண்டிச்சாலை}} (\emph{ciṟṟuṇṭiccālai}) is built from \ta{சிறு}, \textquotedblleft{}small\textquotedblright{} (DEDR 1594), \ta{உண்டி}, \textquotedblleft{}food\textquotedblright{} (DEDR 600), and \ta{சாலை}, \textquotedblleft{}a hall\textquotedblright{}, borrowed from Sanskrit \emph{śālā} (CDIAL 12414).
\end{cousinweb}

\begin{grammarlens}[title={What you've built}]
One more word for getting around.
\end{grammarlens}

\subsection*{Your turn}
\begin{itemize}
\raggedright
  \item \emph{Say it:} \emph{ciṟṟuṇṭiccālai}
  \item \emph{Say it:} \emph{ciṟṟuṇṭiccālai}, once more
  \item \emph{Say it:} \emph{nilaiyattiṉ ciṟṟuṇṭiccālaiyil vaṭai vāṅkiṉōm.} — We bought vadai at the station canteen.
  \item \emph{Recall:} say the Tamil for a site, then the Tamil for a view, then say \emph{ciṟṟuṇṭiccālai} again
\end{itemize}

\begin{tcolorbox}[breakable,colback=teal!4,colframe=teal!35!black,title={Before you move on}]
What does \ta{சிற்றுண்டிச்சாலை} mean? (\textquotedblleft{}A canteen; a snack bar\textquotedblright{}.) Say it once more.
\end{tcolorbox}
